\usepackage{amsmath, amssymb, amsthm}
\usepackage{geometry}
\geometry{margin=1in}
\usepackage{fancyhdr}
% \usepackage[breaklinks=True]{hyperref}
\usepackage[bookmarks]{hyperref}
\usepackage{tikz}
\usepackage{hyperxmp}
\usepackage{enumitem}
\usetikzlibrary{calc}

\newcommand{\PP}{\mathbb{P}} % polynomials
\newcommand{\CC}{\mathbb{C}} % complex numbers
\newcommand{\RR}{\mathbb{R}} % real numbers
\newcommand{\QQ}{\mathbb{Q}} % rational numbers
\newcommand{\NN}{\mathbb{N}} % nonnegative integers
\newcommand{\calF}{\mathcal{F}}
\newcommand{\Z}[1]{\mathbb{Z}/#1\mathbb{Z}} % integers modulo k
                                            % (syntax: "\Z{k}")
\newcommand{\ZZ}{\mathbb{Z}} % integers
\newcommand{\id}{\operatorname{id}} % identity map
\newcommand{\lcm}{\operatorname{lcm}}
% Lowest common multiple. For historical reasons, LaTeX has a \gcd
% command built in, but not an \lcm command. The preceding line
% rectifies that.
\newcommand{\diag}{\operatorname{diag}} % diagonal matrix
\newcommand{\ord}{\operatorname{ord}} % order of a group element
\newcommand{\rank}{\operatorname{rank}} % rank of a matrix
\newcommand{\Hom}{\operatorname{Hom}} % group of homomorphisms
\newcommand{\End}{\operatorname{End}} % ring of endomorphisms
\newcommand{\Iso}{\operatorname{Iso}} % set of isomorphisms
\newcommand{\GL}{\operatorname{GL}} % general linear group
\newcommand{\SL}{\operatorname{SL}} % special linear group
\newcommand{\Nil}{\operatorname{Nil}} % nilradical
\newcommand{\Jac}{\operatorname{Jac}} % Jacobson radical
\newcommand{\Tor}{\operatorname{Tor}} % torsion submodule or Tor functor
\newcommand{\Ext}{\operatorname{Ext}} % Ext functor
\newcommand{\Ker}{\operatorname{Ker}} % kernel
\newcommand{\Coker}{\operatorname{Coker}} % cokernel
\newcommand{\symd}{\mathbin{\bigtriangleup}} % symmetric difference of sets
\newcommand{\powset}[2][]{\ifthenelse{\equal{#2}{}}{\mathcal{P}\left(#1\right)}{\mathcal{P}_{#1}\left(#2\right)}}
% $\powset[k]{S}$ stands for the set of all $k$-element subsets of
% $S$. The argument $k$ is optional, and if not provided, the result
% is the whole powerset of $S$.
\newcommand{\set}[1]{\left\{ #1 \right\}}
% $\set{...}$ compiles to {...} (set-brackets).
\newcommand{\seq}[1]{\left\{ #1 \right\}}
% $\seq{...}$ compiles to {...} (seq-brackets).
\newcommand{\abs}[1]{\left| #1 \right|}
% $\abs{...}$ compiles to |...| (absolute value, or size of a set).
\newcommand{\tup}[1]{\left( #1 \right)}
% $\tup{...}$ compiles to (...) (parentheses, or tuple-brackets).
\newcommand{\ive}[1]{\left[ #1 \right]}
% $\ive{...}$ compiles to [...] (Iverson bracket, aka truth value).
\newcommand{\floor}[1]{\left\lfloor #1 \right\rfloor}
% $\floor{...}$ compiles to |_..._| (floor function).
\newcommand{\underbrack}[2]{\underbrace{#1}_{\substack{#2}}}
% $\underbrack{...1}{...2}$ yields
% $\underbrace{...1}_{\substack{...2}}$. This is useful for doing
% local rewriting transformations on mathematical expressions with
% justifications. For example, try this out:
% $ \underbrack{(a+b)^2}{= a^2 + 2ab + b^2 \\ \text{(by the binomial formula)}} $
\newcommand{\horrule}[1]{\rule{\linewidth}{#1}} % Create horizontal rule command with 1 argument of height
\newcommand{\nnn}{\nonumber\\} % Don't number this line in an "align" environment, and move on to the next line.
\newcommand{\explain}[1]{\tup{\hspace{-0.5pc}\text{\begin{tabular}{c}#1\end{tabular}\hspace{-0.5pc}}}}
% $\explain{...}$ interprets the "..." as text (you can
% nest some math inside by putting it in $...$'s as you
% would in any other text block) and puts it in parentheses.

% Arrows for normal use:
\newcommand{\mono}{\hookrightarrow} % Monomorphism/injection.
\newcommand{\epi}{\twoheadrightarrow} % Epimorphism/surjection.
\newcommand{\iso}{\overset{\cong}{\to}} % Epimorphism/surjection.

% Arrows for xymatrix:
\newcommand{\arinj}{\ar@{_{(}->}} % Monomorphism.
\newcommand{\arinjrev}{\ar@{^{(}->}} % Monomorphism.
\newcommand{\arsurj}{\ar@{->>}} % Epimorphism.
\newcommand{\arelem}{\ar@{|->}} % Arrow to be used in describing a map on elements.
\newcommand{\arback}{\ar@{<-}} % Backward arrow.

%----------------------------------------------------------------------------------------
%	MAKING SUMMATION SIGNS ALWAYS PUT THEIR BOUNDS ABOVE AND BELOW
%	THE SIGN
%----------------------------------------------------------------------------------------
% The following are hacks to ensure that sums (such as
% $\sum_{k=1}^n k$) always put their bounds (i.e., the $k=1$ and the
% $n$) underneath and above the sign, as opposed to on its right.
% Same for products (\prod), set unions (\bigcup) and set
% intersections (\bigcap). Remove the 8 lines below if you do not want
% this behavior.
\let\sumnonlimits\sum
\let\prodnonlimits\prod
\let\cupnonlimits\bigcup
\let\capnonlimits\bigcap
\let\maxnonlimits\max
\let\minnonlimits\min
\let\oldlimits\lim
\renewcommand{\sum}{\sumnonlimits\limits}
\renewcommand{\prod}{\prodnonlimits\limits}
\renewcommand{\bigcup}{\cupnonlimits\limits}
\renewcommand{\bigcap}{\capnonlimits\limits}
\renewcommand{\max}{\maxnonlimits\limits}
\renewcommand{\min}{\minnonlimits\limits}
\renewcommand{\lim}{\oldlimits\limits}

%----------------------------------------------------------------------------------------
%	ENVIRONMENTS
%----------------------------------------------------------------------------------------
% The incantations below define how theorem environments
% (\begin{theorem} ... \end{theorem}) and their likes will look like.
\newtheoremstyle{plainsl}% <name>
  {8pt plus 2pt minus 4pt}% <Space above>
  {8pt plus 2pt minus 4pt}% <Space below>
  {\slshape}% <Body font>
  {0pt}% <Indent amount>
  {\bfseries}% <Theorem head font>
  {.}% <Punctuation after theorem head>
  {5pt plus 1pt minus 1pt}% <Space after theorem headi>
  {}% <Theorem head spec (can be left empty, meaning `normal')>

% Environments which make the text inside them slanted:
\theoremstyle{plainsl}
  \newtheorem{theorem}{Theorem}[section]
  \newtheorem{proposition}[theorem]{Proposition}
  \newtheorem{lemma}[theorem]{Lemma}
  \newtheorem{corollary}[theorem]{Corollary}
  \newtheorem{conjecture}[theorem]{Conjecture}
% Environments that don't:
\theoremstyle{definition}
  \newtheorem{definition}[theorem]{Definition}
  \newtheorem{example}[theorem]{Example}
  \newtheorem{exercise}[theorem]{Exercise}
  \newtheorem{examples}[theorem]{Examples}
  \newtheorem{algorithm}[theorem]{Algorithm}
  \newtheorem{question}[theorem]{Question}
 \theoremstyle{remark}
  \newtheorem{remark}[theorem]{Remark}
\newenvironment{statement}{\begin{quote}}{\end{quote}}
\newenvironment{fineprint}{\begin{small}}{\end{small}}

\newcommand{\sol}{\subsection*{Solution}}
\newcommand{\prob}[1]{\section*{Problem #1}}
\newcommand{\letterbullet}[1]{\item[\textbf{(#1)}]}
\newcommand{\compsign}{(*)}
\newcommand{\diff}[2]{\frac{d#1}{d#2}}
\newcommand{\ydiff}[1]{\diff{y}{#1}}
\newcommand{\pdiff}[2]{\frac{\partial #1}{\partial #2}}
\newcommand{\eval}[1]{\left. #1 \right\rvert}
\newcommand{\st}{\;:\;}
\newcommand{\grad}{\nabla}
\newcommand{\inv}[1]{#1^{-1}}


\newcommand{\eps}{\varepsilon}

\newcommand{\rref}[1]{\xrightarrow{#1}}
\newcommand{\interchange}{\leftrightarrow}

\newcommand{\inner}[2]{\langle #1, #2 \rangle}

\newcommand{\norm}[1]{\left\| #1 \right\|}

\newcommand{\zero}[1]{\mathbf{0}_{#1}}

\newcommand{\Real}[1]{\operatorname{Re}(#1)}

% the for all symbol pisses me off without a space after it. this fixes it
\let\oldforall\forall
\renewcommand{\forall}{\oldforall\;}

% the exists symbol pisses me off without a space after it. this fixes it
\let\oldexists\exists
\renewcommand{\exists}{\oldexists\;}

\newcommand{\recall}{\textbf{Recall:} }

\DeclareMathOperator{\proj}{proj}